\section{Study design}
\label{sec:design}

\subsection{Pre-registration}

The design was written down, committed and hash-stamped on 2026-08-25,
before any model was fit to real HMD data; its SHA-256 digest is
\begin{center}
\small\texttt{bbbe3860a5446d063e186e56555afa893b07c36cf15a990d07567471affe50be}
\end{center}
The registration fixes the analytic choices: regimes and their exact years,
population lists, the model families and mechanisms, the admissible grid,
every metric and its convention, the inference procedure, the validation
gates, and the five hypotheses. It does not, and could not, hide the outcome
of 2020--2024 from us; pre-registration here protects against analytic
flexibility, not against hypothesis surprise. One ordering deviation is
recorded honestly: the registration was frozen before
\citet{dowd2010backtesting} had been read. Nothing in the frozen design
depends on that paper's content; what it can change is the framing of
novelty in the text, which is not frozen. One addendum has been registered
(Addendum~1, 2026-08-26, still before any real-data fit): it adds the
placebo strata and the denominator sensitivities described in
Section~\ref{sec:data}, and changes nothing in the hypotheses, populations,
splits, horizons, metrics or inference plan. \todo{cite the commit hash of
Addendum~1} Every subsequent deviation is reported as a deviation in
Section~\ref{sec:design-deviations}.

\subsection{Three regimes}
\label{sec:design-regimes}

\begin{table}[htbp]\centering
\caption{The three pre-registered evaluation regimes.}
\label{tab:regimes}
\small
\begin{tabular}{@{}lllcc@{}}
\toprule
Regime & Training years & Test years & $h$ & Populations \\
\midrule
Stable (control) & $\le T$, $T=1990,1992,\dots,2014$ & $T{+}1,\dots,T{+}5$ ($\le 2019$) & $1$--$5$ & 20 \\
Shift (primary)  & $\le 2019$ & 2020--2024 & $1$--$5$ & 20 \\
Placebo break    & $\le 1913$ & 1914--1922 & $1$--$9$ & 11 \\
\bottomrule
\end{tabular}
\end{table}

The \emph{stable} regime measures calibration where nothing breaks; it is
the reference level against which \Htwo\ and \Hfive\ are stated. The
\emph{shift} regime is the audit. The \emph{placebo} regime is a second
event---the First World War and the 1918 influenza pandemic---with a longer
horizon, used descriptively: a side-by-side of 1914--1922 and 2020--2024
coverage answers ``is this just COVID?'' without a transfer regression.
Every window is expanding-origin; training always uses all years up to the
origin from the population's own first year. A random split is never used
anywhere, because mortality panels leak trivially under random splitting:
adjacent ages and years are near-duplicates.

\subsection{Model families}
\label{sec:design-families}

Ten families, chosen to span the classical, statistical-learning and neural
lineages of Section~\ref{sec:related}. Each is specified by its reference
paper; the repository documents any deviation.

\begin{enumerate}
\item \textbf{Lee--Carter (SVD).} $\log m_{x,t}=\alpha_x+\beta_x\kappa_t$,
  SVD fit with $\sum_x\beta_x=1$, $\sum_t\kappa_t=0$; random walk with drift
  on $\kappa_t$ including drift-estimation uncertainty
  \citep{lee1992modeling}.
\item \textbf{Poisson Lee--Carter.} Same structure, Poisson likelihood with
  log-exposure offset, fit by Newton's method \citep{brouhns2002poisson}.
\item \textbf{CBD (M5).} $\operatorname{logit} q_{x,t} = \kappa^{(1)}_t +
  \kappa^{(2)}_t (x-\bar x)$ on ages 55--99, bivariate random walk with drift
  \citep{cairns2006two}.
\item \textbf{APC / Renshaw--Haberman (M2-A).} Lee--Carter with a cohort term
  $\gamma_{t-x}$ under the identifiability constraints of
  \citet{cairns2009quantitative}; ARIMA projection of the cohort index.
  \citep{renshaw2006cohort}.
\item \textbf{Sparse VAR.} Banded vector autoregression on log-rate
  improvements with a penalised fit \citep{li2017coherent}.
\item \textbf{Multi-output Gaussian process.} GP regression over (age, year)
  with a shared-across-populations kernel structure
  \citep{huynh2021multi}. \todo{final kernel and inducing-point choices;
  document exclusion if the GPU budget forbids}
\item \textbf{Neural-LC.} Country and sex embeddings feeding a network that
  generalises the Lee--Carter surface \citep{richman2021neural}.
\item \textbf{Shallow CNN-LC.} Convolutional generalisation of Lee--Carter
  \citep{perla2021time}.
\item \textbf{LSTM-on-$\kappa_t$.} Lee--Carter age structure with a recurrent
  network in place of the random walk \citep{nigri2019deep}.
\item \textbf{Distributional Poisson/NB head.} A network whose output layer
  parameterises a Poisson or negative-binomial count distribution for
  $D_{x,t}$ given $\log E_{x,t}$; the only neural family with a model-native
  predictive distribution.
\end{enumerate}

Hyperparameters of the neural families are tuned on an inner \emph{time}
split---the final five pre-cutoff years of each training window---never on a
test year and never on a random subset.

\subsection{Uncertainty mechanisms}
\label{sec:design-mechanisms}

Seven mechanisms, each a way of turning a fitted family into predictive
samples.

\begin{enumerate}
\item \textbf{Model-native.} The predictive distribution the reference paper
  itself supplies: random-walk innovation and drift uncertainty for
  Lee--Carter-type models, the GP posterior, the count distribution of the
  distributional head.
\item \textbf{Semiparametric Poisson bootstrap.} Refit the family $B$ times to
  pseudo-deaths $D^{(b)}\sim\mathrm{Poisson}(E\hat m)$ and pool the refits'
  native samples \citep{brouhns2005bootstrapping}; adds parameter uncertainty
  the native mechanism conditions away.
\item \textbf{Deep ensemble.} $M=10$ independently seeded trainings, samples
  pooled with equal weight \citep{lakshminarayanan2017simple}.
\item \textbf{Monte Carlo dropout.} Dropout active at prediction time,
  stochastic forward passes as samples \citep{gal2016dropout}.
\item \textbf{Split conformal.} Absolute-residual conformal intervals on the
  log-rate scale \citep{lei2018distribution}, Mondrian-stratified by age
  band ($0$--$24$, $25$--$64$, $65$--$99$) and by horizon
  \citep{vovk2005algorithmic}.
\item \textbf{EnbPI.} Ensemble batch prediction intervals adapted from the
  online original to annual multi-step forecasting \citep{xu2021conformal};
  the adaptation is documented in the repository.
\item \textbf{Copula / joint-path conformal.} A simultaneous band over
  horizons from a scaled sup-norm nonconformity score
  \citep{diquigiovanni2022conformal}, in the spirit of
  \citet{sun2024copula} and \citet{messoudi2021copula}; the only mechanism
  whose band is calibrated to be joint over $h=1,\dots,H$.
\end{enumerate}

Two disciplines apply to every conformal arm. The calibration residuals are
drawn only from years at or before the training cutoff, on an inner time
split in which the calibration years are later than the years the base fit
is trained on; no mechanism sees a test year. And a conformal wrapper yields
an interval, not a distribution: to preserve the single sample interface the
harness draws uniformly inside each cell's interval, independently across
cells. That convention keeps coverage, width and interval score exact at the
construction level, but it makes CRPS, log score and PIT for conformal cells
a placeholder rather than a distributional claim. Those scores are computed,
flagged, and reported only in Appendix~\ref{app:secondary}; they are never
ranked against distributional mechanisms.

\subsection{The crossed grid}
\label{sec:design-grid}

Not every family admits every mechanism. Deterministic classical fits have
no seed variance, so ensembles and dropout are undefined for them; the GP
posterior already integrates parameter uncertainty, so a bootstrap would be
redundant; and neural networks have no closed-form native predictive except
through the distributional head. Table~\ref{tab:grid} enumerates the
admissible cells: 50 primary and a handful of secondary ones. The three
conformal columns are wrappers around the same underlying fits as the native
and ensemble columns, not refits, so the expensive axis is neural families
$\times$ ensemble members $\times$ expanding origins.

\begin{table}[htbp]\centering
\caption{Admissible model-family $\times$ uncertainty-mechanism grid, transcribed
from the pre-registration.}
\label{tab:grid}
\inputtable{tab-grid}
\end{table}

Because the grid is ragged, no claim in this paper is a full-factorial main
effect. Every comparison is a \emph{contrast within an admissible sub-grid},
and the pre-registration names the contrasts that carry the hypotheses:
conformal versus native within the classical families; ensemble versus
dropout within the neural families; bootstrap versus native within the
classical families; and joint-path conformal versus split conformal across
all families.

\subsection{One interface, one scoring path}
\label{sec:design-interface}

Every family--mechanism cell implements the same two calls,
\[
  \texttt{fit}(D, E) \quad\text{and}\quad
  \texttt{sample\_mx}(h, n, \text{rng}) \;\to\; [n, h, n_{\text{ages}}],
\]
returning $n$ pathwise trajectories of $m_{x,T+1},\dots,m_{x,T+h}$. Point
forecasts, intervals, proper scores, PIT values, joint path coverage and the
life-table functionals are all computed from these samples by one evaluation
module. No model has a bespoke metric implementation, so no model can be
accidentally flattered by one. The point forecast for RMSE and MAE is the
pointwise median of the log-rate samples. \todo{unify the point functional
for $e_0$ (currently the sample mean) with the median used for log rates}

\subsection{Seeds}

The global seed is 20260825. Each cell derives its own generator from a seed
sequence over (global seed, origin, population, sex, family, mechanism);
ensemble member $j$ uses global seed plus $j$. Every seed is written to the
results file alongside the metrics.

\subsection{Validation gates}
\label{sec:design-gates}

Three gates had to pass before any real-data result could be read. All three
have passed.

\paragraph{Gate 1: synthetic truth.} Data are simulated from a known Poisson
Lee--Carter process (40 ages, 60 training years, horizons $1$--$10$,
exposures $10^5$, thirty independent worlds, 1{,}500 predictive samples
each). The correctly specified Poisson Lee--Carter must attain nominal
95\% coverage within Monte Carlo tolerance and a near-uniform PIT histogram;
the harness must also detect an injected break as under-coverage. It does
both. If this gate had failed, every coverage number on real data would have
been measuring a bug. One documented finding from the gate: the two-stage
SVD Lee--Carter is mildly \emph{over}-dispersed on Poisson data, because the
SVD $\kappa_t$ carries observation noise into the random-walk innovation
variance; its intervals are slightly wide and its PIT centre-heavy. This is
a property of the method, not of the harness, and it is carried into the
interpretation of the SVD-LC rows.

\paragraph{Gate 2: oracle parity.} The Python Poisson Lee--Carter must match
R \pkg{StMoMo} (version 0.4.1 under R~4.6.1) on a shared subset. On Swedish
males, 1950--2000, ages $0$--$99$, the two fits have identical Poisson
log-likelihood and a maximum relative difference in $\beta_x$ of
$2.4\times10^{-14}$ when the Newton iteration is run to its cap.
\todo{quote the $\alpha_x$ and $\kappa_t$ relative differences from the
parity script} The \pkg{StMoMo} parameter dump is
tracked in the repository so the check can be re-run without R.

\paragraph{Gate 3: life-table parity.} Our period life table on HMD's
published $m_x$ reproduces HMD's published $e_0$ and $e_{65}$ within
0.15~years at spot checks (Section~\ref{sec:data}).

\subsection{Inference plan}

The shift regime is one event. Twenty populations experienced one common
shock, so the effective number of independent observations of ``what a
pandemic does to an interval'' is nearer one than twenty, and nearer one
than the thousands of cells in the panel. Inference is therefore framed as
an event study: population is the cluster, every test resamples clusters
rather than cells, PIT is pooled within cluster before it is compared
across clusters, and the placebo regime is treated as a second event rather
than as more draws from the same one. The procedures are defined in
Section~\ref{sec:metrics-inference}.

\subsection{Deviations from the pre-registration}
\label{sec:design-deviations}

\begin{itemize}
\item Ordering: \citet{dowd2010backtesting} read after the registration was
  frozen (above).
\item Addendum~1 (2026-08-26): pre-specified placebo strata and shift-regime
  denominator sensitivities. Not a deviation---nothing frozen changed---but
  recorded here so the ledger is complete.
\item \todo{append every subsequent deviation here with its date and reason;
  an empty list at submission is a claim, not a default}
\end{itemize}
